\section{Discussion}
\label{sec:discussion}

The main result is not that a classifier can recognize a simulated failure
trajectory after seeing its age. Cycle count alone already reaches AUROC 0.923.
The useful result is the comparison against that strong baseline. At a matched
pre-window alert budget, measured telemetry raises final-window recall from
0.554 for cycle plus operating conditions to 0.793 for all observable inputs.
The paired complete-engine bootstrap interval for the recall gain remains
positive, from 0.183 to 0.296.

The secondary nested model-selection policy raises all-observable recall from
0.793 to 0.821 and AUROC from 0.966 to 0.973. It selects histogram gradient
boosting in three outer folds and an RBF support-vector classifier in two. This
is a practical improvement under the same warning budget, but it is not an
architecture claim. The candidate set is small, the model choice is evaluated
only on this simulated subset, and the policy does not improve the earlier
first-warning timing reported for the sensor-only rule.

The telemetry-only rule supplies a separate result for an early-screening
setting. It has no elapsed-cycle feature, yet it recalls 77.1\% of final-window
cycles at a 4.7\% pre-window false-alert rate. It warns 69 of 90 engines by
the window start, with a median lead of 6 cycles among warned engines. This
does not show that the model knows when a real engine should be removed from
service. It shows that the documented measured channels contain signal about
the simulated approach to the chosen 20-cycle boundary.

The feature definition matters. The evaluated model uses only the 14 named
measured physical channels, listed in Table~\ref{tab:feature-provenance}.
It does not use virtual sensors, target health parameters, unit ID, RUL, or the
auxiliary \texttt{Fc} and \texttt{hs} fields. This restriction makes the
telemetry result testable. It also makes the result lower than an experiment
that treated auxiliary health-state information as if it were a measured
sensor.

Several limitations remain. First, tiny-N-CMAPSS is simulated and the retained
data are $100\times$ downsampled. Each per-cycle summary averages about 93
retained readings, so it cannot test how a warning would respond to high-rate
within-flight transients. Second, the 20-cycle window, 5\% false-alert budget,
and one-cycle alert definition are experimental choices. A maintenance program
would need costs, inspection capacity, confirmation rules, and a lead-time
objective before choosing a policy. Third, five seeds test stochastic variation
of the fixed reference and the nested selection procedure on the same data.
They do not replace an independently labeled fleet or a held-out test campaign.

The next technical step is to train and evaluate a warning policy against an
engine-level objective, such as warning by a specified lead time with a
per-engine alert budget. That would directly optimize the tradeoff exposed in
Table~\ref{tab:warning-timing}. The next empirical step is external validation
on labeled engines with maintenance outcomes. Until then, this paper should be
read as a controlled simulation result about information value, not a deployable
maintenance recommendation.
